\section{Introduction}
\label{sec:introduction}

\noindent
This paper is organized in two parts: chapter~\ref{chap:1}~goes over the preliminaries and theoretical background of gravitational waves, while chapter~\ref{chap:2}~delves into the Einstein field equations, linearized gravity, and recent developments in gravitational wave astronomy. We don't assume any prior knowledge (except the fundamentals, like calculus, linear algebra, etc), as we develop those throughout the course of the paper.

If you already have a strong understanding of differential geometry, you can skip to section~\ref{sec:theoretical_physics_background}, where we introduce the necessary physics background for understanding gravitational waves.

Please note that this paper will be more mathematical in nature, as I aim to explore the mathematical framework underlying gravitational waves, rather than focusing solely on their physical implications. The mathematical perspective will allow us to gain a deeper understanding of the properties and behavior of gravitational waves, as well as their significance in the broader context of physics and astronomy.

The goal is to provide a comprehensive overview of the mathematical framework underlying gravitational waves, while also highlighting their significance in the broader context of physics and astronomy. We start this chapter by going over the notation used throughout this paper.
